\chapter{Arbitrage de crédit relatif}\label{ch:s2:credit-relative-value}

L'obligation d'une entreprise et son credit default swap assurent le même défaut. Acheter l'obligation et acheter la protection ne devraient rapporter que le
taux sans risque, de sorte que la base, soit le spread du CDS moins celui de l'obligation, devrait être nulle. Bai et Collin-Dufresne constatent que l'écart est le plus marqué pour les obligations soumises aux
frictions les plus fortes~: peu liquides, coûteuses à financer, exposées à la contrepartie, mauvaises en tant que collatéral. En 2008, la base a fortement basculé en territoire négatif lorsque la transaction
censée la refermer exigeait un financement dont personne ne disposait. Sur le marché synthétique de ce chapitre, la base se situe à $-15$ points de base jusqu'à ce qu'une crise de financement
la porte à $-151$. Un book de base négative à effet de levier dix fois perd 60~\% de son capital dans la baisse et, s'il survit, se redresse de 76~\% après le point bas.
Entré au point bas, le même trade rapporte 70~\% en un an. La construction repose sur \texttt{firm.creditrv}.

\section{La base obligation-CDS}

Le spread d'une obligation (livre~2, chapitre~21, le Z-spread par rapport à la courbe sans risque) et un credit default swap sur le même émetteur (livre~2, chapitre~23) valorisent tous deux
le défaut de l'émetteur. Une obligation achetée et couverte par une protection CDS ne porte plus de risque de défaut~: en cas d'événement de crédit, la protection verse le pair diminué
du taux de recouvrement de l'obligation. Sur des marchés sans frictions, le spread obligataire est égal au spread du CDS. En pratique, l'obligation doit être financée, le CDS exige un collatéral, et les
deux côtés sont négociés par des investisseurs différents.

\begin{definition}[Trade de base négative]\label{def:s2:credit-relative-value:negbasis}
Un \emph{trade de base négative}\index{trade de base négative} achète une obligation dont le spread dépasse celui du CDS de son émetteur, la finance en pension livrée et achète
une protection CDS sur le même émetteur~; il rapporte l'écart entre les deux spreads diminué de son coût de financement, ne porte aucun risque de défaut, et perd sur ses valorisations lorsque l'écart
s'élargit encore.
\end{definition}

\texttt{firm.creditrv} construit quarante émetteurs synthétiques dont les spreads de CDS avoisinent 120 points de base, avec une moyenne de 136 sur l'échantillon. Le spread de chaque obligation est égal à son
spread de CDS, auquel s'ajoute une prime de financement commune à toutes les obligations (15 points de base en temps normal), plus un bruit. La base vaut donc l'opposé de la prime de financement.
À l'année~5, une crise intégrée au modèle porte la prime à 150 points de base sur un trimestre et la ramène à son niveau en un an
(voir la \cref{fig:s2:credit-relative-value:basis}). La base moyenne passe de $-14.9$ à $-150.7$ points de base, à l'année~5.25.

\section{Indice contre composantes}

\begin{definition}[Arbitrage d'indices de crédit]\label{def:s2:credit-relative-value:index}
L'\emph{arbitrage d'indices de crédit}\index{arbitrage d'indices de crédit} négocie un indice de CDS contre le panier des CDS mono-nom de ses composantes lorsque le spread de l'indice diffère de la
moyenne des spreads des membres (le skew d'indice)~: on vend la protection du côté le moins cher et on l'achète du côté le plus cher.
\end{definition}

L'indice se négocie davantage que ses membres et réagit en premier. L'écart entre les deux, le skew, reflète cette liquidité, et il se referme au fur et à mesure que les noms rattrapent l'indice.
Le skew synthétique est un bruit à retour à la moyenne de 4 points de base, de demi-vie 20~jours. Le négocier dès qu'il dépasse un écart type, puis conserver
la position jusqu'à ce qu'il croise zéro, a rapporté 55 points de base de notionnel par an, avec un ratio de Sharpe de 0.97. Le book a détenu une position 60~\% du temps et a payé 1
point de base de spread (4.5 sur le notionnel) pour négocier chaque unité, car la jambe mono-nom coûte cher.

\section{Trades de courbe de crédit}

\begin{definition}[Trade de courbe de crédit]\label{def:s2:credit-relative-value:curve}
Un \emph{trade de courbe de crédit}\index{trade de courbe de crédit} achète la protection à une échéance et la vend à une autre, sur le même émetteur ou indice, en montants qui compensent leurs
durations de spread, en pariant sur la pente de la courbe de crédit plutôt que sur son niveau.
\end{definition}

Un flattener vend la protection à l'échéance courte et l'achète à l'échéance longue, et gagne lorsque le spread long baisse par rapport au spread court. Il porte
un risque de défaut entre les deux échéances~: un défaut survenant avant l'échéance du contrat court donne lieu au paiement des deux jambes, mais un défaut survenant après ne donne lieu qu'au paiement de la jambe longue. La pente synthétique
dix ans moins cinq ans affiche une moyenne de 40 points de base et revient vers sa moyenne avec une demi-vie de 60~jours. Le trade en bande a rapporté 49 points de base par an, avec un ratio de Sharpe
de 0.75, pour une position détenue 62~\% du temps.

\section{Financement et base négative}

Le book de base négative détient chaque obligation contre une protection sur son émetteur, finance les obligations avec une décote de 10~\%, et paie son propre spread de financement sur
la part financée (voir le \cref{lst:s2:credit-relative-value:book}). Le capital correspond à la décote, de sorte que le book est à effet de levier dix fois.

\begin{center}
\small
\setlength{\tabcolsep}{4pt}%
\begin{tabular}{@{}lcccc@{}}
\toprule
sur capital, dix fois & avant (\%/an) & jusqu'au point bas (\%) & après le point bas (\%) & dix ans (\%)\\
\midrule
financement propre 0 pb & 1.6 & $-59.9$ & 76.0 & 24.4\\
financement propre 30 pb & $-1.1$ & $-60.6$ & 63.2 & $-2.6$\\
financement propre 60 pb & $-3.8$ & $-61.2$ & 50.4 & $-29.6$\\
\bottomrule
\end{tabular}
\end{center}

Avant la crise, le trade rapporte à peine. Une base de $-15$ points de base couvre mal le financement, et un trader qui se finance à 30 points de base perd
de l'argent en le détenant. Lorsque la crise élargit la base de 135 points de base, une duration de spread de 4.5 transforme cet élargissement en une perte de 6.1~\% sur le notionnel, soit 61~\% du
capital à dix fois (voir la \cref{fig:s2:credit-relative-value:basis}). Aucun défaut n'est survenu et la couverture a fonctionné. La perte est une perte de valorisation, et un book qui doit faire face
à des appels de marge vend au point bas sans jamais voir le redressement. Mitchell et Pulvino ont décrit ce mécanisme en 2008~: la défaillance imminente des
prime brokers a réduit brutalement l'effet de levier des hedge funds, et les arbitragistes qui auraient dû acheter des obligations bon marché ont dû les vendre.

L'autre face de la crise est le meilleur trade du book. Entré au point bas de la base et conservé un an, le book de base négative a rapporté 69.8~\% de son
capital sans coût de financement, et 67.1~\% avec un financement à 30 points de base. Le carry en a fourni 8.4 et 5.7 points, et le retour de la base à la normale
les 61.4 points restants. Le capital capable d'acheter au point bas est celui qui n'était pas poussé à la limite de son levier avant la crise.

\begin{omfigure}
\begin{tikzpicture}
\begin{axis}[width=11cm,height=5.4cm,xlabel={\small année},ylabel={\small pb ou~\% du capital},tick label style={font=\footnotesize},
  xmin=0,xmax=10,ymin=-160,ymax=40,grid=major,grid style={gray!25},table/col sep=comma,
  legend style={font=\scriptsize,at={(0.5,-0.3)},anchor=north},legend columns=3]
\addplot[black,thick] table[x=year,y=basis]{figdata/strategies-2/18-credit-relative-value/basis.csv};
\addplot[omThm,thick] table[x=year,y=book0]{figdata/strategies-2/18-credit-relative-value/basis.csv};
\addplot[omDef,thick,dashed] table[x=year,y=book30]{figdata/strategies-2/18-credit-relative-value/basis.csv};
\legend{base moyenne (pb),book à financement 0 pb (\%),book à financement 30 pb (\%)}
\end{axis}
\end{tikzpicture}

\omcaption{La base obligation-CDS synthétique sur dix ans, avec une crise de financement en année~5, et le rendement cumulé sur capital d'un book de base négative
détenu depuis le début, à effet de levier dix fois, pour deux coûts de financement. Données~: \texttt{s2\_creditrv.market}.}\label{fig:s2:credit-relative-value:basis}
\end{omfigure}

\section{Fiches de stratégie}

\begin{strategyfile}{Trade de base négative}\label{strat:s2:credit-relative-value:negbasis}
\sfield{Qui paie, et pourquoi} Les détenteurs contraints de vendre des obligations lorsque le financement se raréfie~; la base rémunère le bilan qui les porte.
\sfield{Instruments et plateformes} Obligations d'entreprises en pension livrée~; CDS mono-nom.
\sfield{Signal} La base comparée à son propre financement et à sa décote.
\sfield{Dimensionnement et exécution} Effet de levier calibré sur un élargissement de crise et non sur la base normale~; financement à terme.
\sfield{Coûts} Financement, décotes, spreads de CDS et collatéral.
\sfield{Comment elle échoue} Une crise de financement qui élargit la base pendant que le levier est retiré.
\sfield{Horizon, capacité, infrastructure} De quelques mois à quelques années~; lignes de repo et documentation des CDS.
\sfield{Backtester honnêtement} Financement et décotes tels qu'ils ont évolué~; ventes forcées au point bas.
\sfield{Sources} Bai et Collin-Dufresne (2019)~; Mitchell et Pulvino (2012)~; ce chapitre~: $-60~\%$ du capital à dix fois pendant la crise, $+70~\%$ dans
l'année qui suit son point bas.
\end{strategyfile}

\begin{strategyfile}{Arbitrage d'indice contre arbitrage mono-nom}\label{strat:s2:credit-relative-value:index}
\sfield{Qui paie, et pourquoi} Traders d'indices qui font bouger l'indice avant les noms individuels.
\sfield{Instruments et plateformes} Indices de CDS~; les CDS mono-nom des membres.
\sfield{Signal} Le skew d'indice comparé à son historique.
\sfield{Dimensionnement et exécution} Entrée à bandes, sortie à la convergence.
\sfield{Coûts} Les spreads de CDS mono-nom, le coût le plus élevé.
\sfield{Comment elle échoue} Noms peu liquides~; événements de crédit chez les membres entre deux rolls.
\sfield{Horizon, capacité, infrastructure} Quelques semaines~; un système d'exécution de panier.
\sfield{Backtester honnêtement} Coûts des CDS mono-nom sur l'ensemble du panier.
\sfield{Sources} Ce chapitre~: 55~points de base par an, ratio de Sharpe 0.97 (skew injecté).
\end{strategyfile}

\begin{strategyfile}{Flattener de courbe de crédit}\label{strat:s2:credit-relative-value:curve}
\sfield{Qui paie, et pourquoi} Demande de protection à une échéance donnée (couverture courte en période de stress, longue issue des produits structurés).
\sfield{Instruments et plateformes} CDS à deux échéances, sur noms ou sur indices.
\sfield{Signal} La pente comparée à son historique et à la qualité de crédit de l'émetteur.
\sfield{Dimensionnement et exécution} Duration de spread appariée.
\sfield{Coûts} Deux jambes.
\sfield{Comment elle échoue} Défaut entre les deux échéances~; courbe qui s'inverse en période de détresse.
\sfield{Horizon, capacité, infrastructure} Quelques mois.
\sfield{Backtester honnêtement} Saut au défaut sur la jambe mal appariée.
\sfield{Sources} Ce chapitre~: 49~points de base par an, ratio de Sharpe 0.75 (pente injectée).
\end{strategyfile}

\begin{strategyfile}{Dette senior contre dette subordonnée}\label{strat:s2:credit-relative-value:seniority}
\sfield{Qui paie, et pourquoi} Mauvaise évaluation entre des créances de rang de séniorité différent sur un même émetteur.
\sfield{Instruments et plateformes} Obligations senior et subordonnées, ou CDS.
\sfield{Signal} Le rapport de leurs spreads, confronté aux hypothèses de recouvrement implicites.
\sfield{Dimensionnement et exécution} Couvert contre le risque de défaut de l'émetteur.
\sfield{Coûts} Titres subordonnés peu liquides.
\sfield{Comment elle échoue} Règles de renflouement interne (bail-in) ou de restructuration qui modifient les recouvrements.
\sfield{Horizon, capacité, infrastructure} Quelques mois.
\sfield{Backtester honnêtement} Recouvrements issus des résolutions réelles.
\sfield{Sources} Aucun chiffre de performance vérifié.
\end{strategyfile}

\begin{strategyfile}{Crédit multidevise}\label{strat:s2:credit-relative-value:currency}
\sfield{Qui paie, et pourquoi} Obligations d'émetteurs valorisées différemment dans deux devises une fois swappées.
\sfield{Instruments et plateformes} Obligations d'un même émetteur dans deux devises~; swaps de devises.
\sfield{Signal} L'écart de spreads après swap.
\sfield{Dimensionnement et exécution} Risque de change et de taux neutralisé par le swap.
\sfield{Coûts} La base de devises (chapitre~14).
\sfield{Comment elle échoue} La base elle-même se déplace.
\sfield{Horizon, capacité, infrastructure} Quelques mois~; lignes de swap.
\sfield{Backtester honnêtement} Coûts des swaps à l'époque.
\sfield{Sources} Aucun chiffre de performance vérifié.
\end{strategyfile}

\section{Tutoriel~: même défaut, à nouveau deux prix}\label{tut:s2:credit-relative-value:tut}

\begin{tutorial}
\textbf{Objectif.} Simuler des obligations et des CDS avec une prime de financement, faire passer le book de base négative au travers d'une crise de financement, et trader le skew d'indice et la
courbe. \textbf{État final~:} le tableau et la figure.
\begin{tutsteps}
\item \textbf{Les books}.
\omcode{code/firm/creditrv/firm_creditrv.py}{84}{111}{Le book de base négative et les trades de relative value par bandes.}\label{lst:s2:credit-relative-value:book}
\item \textbf{La crise}.
\omcode{code/strategies-2/18-credit-relative-value/python/s2_creditrv.py}{46}{66}{Le book traversant la crise, et la position ouverte au point bas.}\label{lst:s2:credit-relative-value:crisis}
\item \textbf{Exécuter} \texttt{basis\_stats()}, \texttt{basis\_book()}, \texttt{at\_the\_trough()}, \texttt{rv\_trades()} et \texttt{fig\_creditrv.py}.
\end{tutsteps}
\textbf{Pour aller plus loin.} Ajouter une règle de marge qui réduit le book lorsque le capital passe sous la décote et mesurer la perte ainsi figée~; rendre la
prime de financement différente selon la liquidité des obligations~; ajouter un défaut.
\end{tutorial}

\section{Construction~: relative value de crédit}\label{bld:s2:credit-relative-value:creditrv}

\begin{build}
\bfield{Rôle} Une base obligation-CDS pilotée par le financement, le book de base négative, les trades sur composants d'indice et de courbe.
\bfield{Interface} \texttt{CreditConfig(\dots)}, \texttt{simulate\_credit(cfg)}, \texttt{negative\_basis(sim, cfg, own\_funding)},
\texttt{index\_arbitrage(sim, cfg)}, \texttt{curve\_trade(sim, cfg)}.
\bfield{Règles} La base est égale à moins la prime de financement plus un bruit~; valorisations à la duration de spread~; transactions par bandes, avec des coûts par unité échangée.
\bfield{Tests d'acceptation} \texttt{code/firm/creditrv/tests/}~: la base est égale à moins la prime de financement sans bruit~; carry et valorisations calculés à la main~; la
règle de position par bandes.
\bfield{Pour aller plus loin} Défauts~; appels de marge~; financement dépendant de la liquidité.
\end{build}

\begin{omsources}
\item J.~Bai and P.~Collin-Dufresne, ``The CDS-bond basis'', \emph{Financial Management} 48(2), 2019.
\item M.~Mitchell and T.~Pulvino, ``Arbitrage crashes and the speed of capital'', \emph{Journal of Financial Economics} 104(3), 2012.
\end{omsources}

\section{Exercices}

\begin{exercise}[$\star$]\label{exo:s2:credit-relative-value:1}
La base vaut $-15$ points de base, un trader se finance à 30~points de base au-dessus du taux sans risque, et la décote est de 10~\%. Quel est le carry annuel sur le capital~?
\end{exercise}

\begin{exercise}[$\star$]\label{exo:s2:credit-relative-value:2}
La base s'élargit de $-15$ à $-150$ points de base avec une duration de spread de 4.5. Quelle perte en capital subit un book avec un effet de levier de dix~?
\end{exercise}

\begin{exercise}[$\star$]\label{exo:s2:credit-relative-value:3}
Quel est le levier maximal qu'autorise une décote de 10~\%~?
\end{exercise}

\begin{exercise}[$\star\star$]\label{exo:s2:credit-relative-value:4}
Pourquoi le \omterm{def:s2:credit-relative-value:negbasis}{trade de base négative} est-il exempt de risque de défaut sans être exempt de perte~?
\end{exercise}

\begin{exercise}[$\star\star$]\label{exo:s2:credit-relative-value:5}
Pourquoi Bai et Collin-Dufresne ont-ils constaté une base plus négative lorsque les frais de prêt de l'obligation sont élevés~?
\end{exercise}

\begin{exercise}[$\star\star$]\label{exo:s2:credit-relative-value:6}
Quel risque de défaut un flattener de courbe de crédit conserve-t-il~?
\end{exercise}

\begin{exercise}[$\star\star\star$]\label{exo:s2:credit-relative-value:7}
\emph{Programmation.} Ajouter une règle qui vend le book de base négative dès que sa perte cumulée atteint 50~\% du capital. Que fige-t-elle, et
ce qu'elle manque~?
\end{exercise}

\begin{exercise}[$\star\star\star$]\label{exo:s2:credit-relative-value:8}
\emph{Trouver la faille.} «~Le basis négatif est un arbitrage~: on le conserve jusqu'à l'échéance et on ne peut pas perdre.~»
\end{exercise}

\section{Problème~: Même défaut, deux prix à nouveau}

\begin{problem}[{Problème du week-end --- valeur relative de crédit}]\label{pb:s2:credit-relative-value:1}
Les émetteurs synthétiques du chapitre et les données publiques.

\medskip\noindent\textbf{Partie I --- La base.}
\begin{enumerate}
\item Définir la base obligation-CDS et le \omterm{def:s2:credit-relative-value:negbasis}{trade de base négative}.
\item Pourquoi la base devrait-elle être nulle, et pourquoi ne l'est-elle pas~?
\item Que Bai et Collin-Dufresne ont-ils constaté~?
\item Décrire la prime de financement synthétique et la crise.
\end{enumerate}

\medskip\noindent\textbf{Partie II --- Autres trades.}
\begin{enumerate}[resume]
\item Définir l'\omterm{def:s2:credit-relative-value:index}{arbitrage d'indices de crédit} et donner son résultat synthétique.
\item Définir un \omterm{def:s2:credit-relative-value:curve}{trade de courbe de crédit} et donner son résultat synthétique.
\item Quels risques ces deux trades portent-ils~?
\item Pourquoi la jambe sur nom unique est-elle la plus coûteuse~?
\end{enumerate}

\medskip\noindent\textbf{Partie III --- La crise.}
\begin{enumerate}[resume]
\item Donner les rendements du book avant la crise pour chaque coût de financement.
\item Donner la perte jusqu'au plus bas et la reprise.
\item Que Mitchell et Pulvino ont-ils décrit~?
\item Que rapportait le trade s'il était entré au plus bas~?
\end{enumerate}

\medskip\noindent\textbf{Partie IV --- Le verdict.}
\begin{enumerate}[resume]
\item Énoncer le \emph{résultat nommé}~: le portage du \omterm{def:s2:credit-relative-value:negbasis}{trade de base négative} et sa perte en valorisation au prix de marché lorsque les coûts de financement bondissent.
\item Pourquoi le trade rapporte-t-il à peine en temps normal~?
\item Qui peut acheter au plus bas~?
\item Comment dimensionneriez-vous un book de base négative~?
\item Comment le backtester honnêtement~?
\item De quelle fiche de stratégie dépend la base de devises~?
\item Quel lien ce chapitre entretient-il avec les chapitres~8 et 11~?
\item En une phrase~: que price la base négative~?
\end{enumerate}
\end{problem}

\section{Questions d'entretien}

\begin{interviewq}[$\star$ \iqroles{trader}]\label{iq:s2:credit-relative-value:1}
Qu'est-ce que la base CDS-obligation, et que signifie une base négative~?
\end{interviewq}

\begin{interviewq}[$\star\star$ \iqroles{trader}]\label{iq:s2:credit-relative-value:2}
Détaillez un \omterm{def:s2:credit-relative-value:negbasis}{trade de base négative} et ses flux de trésorerie, y compris en cas d'événement de crédit.
\end{interviewq}

\begin{interviewq}[$\star\star$ \iqroles{risque}]\label{iq:s2:credit-relative-value:3}
Un book de base négative est endetté à dix fois ses fonds propres. Que stresseriez-vous~?
\end{interviewq}

\begin{interviewq}[$\star\star$ \iqroles{chercheur}]\label{iq:s2:credit-relative-value:4}
Comment calculeriez-vous la valeur intrinsèque d'un indice de CDS, et pourquoi diffère-t-elle de l'indice~?
\end{interviewq}

\begin{interviewq}[$\star\star$ \iqroles{développeur}]\label{iq:s2:credit-relative-value:5}
Quelles données de référence un book de CDS exige-t-il, et qu'est-ce qui déraille lors des roll d'indices~?
\end{interviewq}

\begin{interviewq}[$\star\star\star$ \iqroles{chercheur}]\label{iq:s2:credit-relative-value:6}
Montrez qu'une obligation couverte par une protection CDS jusqu'à l'échéance présente, en négligeant le financement et l'option de livraison,
les flux d'une obligation sans risque plus l'écart de spread, et nommer ce qui rompt l'équivalence en pratique.
\end{interviewq}
